\section{Additional Ablation Measurements}
\label{app:f-ablation}

\noindent\textbf{Acceptance-rule controls.}
The four controls in Table~\ref{tab:f-rule-gating} use call insertion,
$H\!\rightarrow\!B$ boundary tampering, cross-call evidence splicing,
and aggregate-bill manipulation.
For each attack, five fixtures are each evaluated 100 times under the
unmodified decision and the matching relaxed decision.
This gives 500 observations per attack--decision configuration and
4,000 observations overall.
Repeated evaluations of a fixture are not distinct attack constructions.

Let $a(x)$ be the unmodified auditor's Boolean decision and $f(x)$ its
first failed relation for candidate $x$.
The control for relation $R$ computes
\[
 a_R(x)=a(x)\lor\bigl[f(x)=R\bigr].
\]
The auditor executes its normal first-failure logic before this decision
override; the control neither removes a component nor continues checking
later relations after rejection.
Only the four matched attack--rule combinations were measured.
The archived variant named \texttt{w/o-R\_link} implements the
$R_{\mathrm{prov}}$ override and therefore relaxes the entire provider
rejection gate, rather than isolating $R_{\mathrm{link}}$.

\noindent\textbf{Microbenchmark scope.}
Each fixture contains eight calls and 32 lifecycle events.
The five cumulative configurations are the empty request-ID baseline,
lifecycle validation, added bilateral-handoff validation, added provider
validation, and added settlement validation.
For every configuration, each of five runs performs 100 warmups followed
by 1,000 measured executions: 2,500 warmups and 25,000 measured executions
overall.
Evidence and signing keys are prepared before timing.
The timer covers the selected validation operations and construction of
the corresponding payload-field selections.
Authorization verification, evidence generation and signing, model
inference, canonical serialization, persistence, ledger operations, and
proof generation are excluded.

Selected payload size is the canonical encoding of the selected fields
after the timer stops.
It is not the complete persisted evidence bundle.
The baseline contains only a request identifier and occupies 22\,B;
the final selected payload occupies 32,990\,B.
Table~\ref{tab:f-cumulative-cost} reports cumulative wall time,
process CPU time, and selected payload size.

\begin{table}[!t]
\centering
\caption{Cumulative validation costs. Times are mean $\pm$ sample SD
across five run means; payload sizes are identical across runs.}
\label{tab:f-cumulative-cost}
\footnotesize
\renewcommand{\arraystretch}{1.08}
\setlength{\tabcolsep}{3pt}
\begin{tabular*}{\columnwidth}{@{\hspace{3pt}}l@{\extracolsep{\fill}}ccc@{\hspace{3pt}}}
\toprule
\shortstack{\textbf{Cumulative}\\\textbf{checks}}
& \shortstack{\textbf{Wall time}\\\textbf{(ms)}}
& \shortstack{\textbf{CPU time}\\\textbf{(ms)}}
& \shortstack{\textbf{Payload}\\\textbf{(KiB)}} \\
\midrule
Baseline & 0.00047 $\pm$ 0.00007 & 0.00129 $\pm$ 0.00022 & 0.0215 \\
+ Lifecycle & 6.026 $\pm$ 0.029 & 6.020 $\pm$ 0.031 & 15.369 \\
+ Handoff & 11.452 $\pm$ 0.037 & 11.432 $\pm$ 0.030 & 22.468 \\
+ Provider & 14.361 $\pm$ 0.095 & 14.342 $\pm$ 0.092 & 29.543 \\
+ Settlement & 16.704 $\pm$ 0.109 & 16.673 $\pm$ 0.103 & 32.217 \\
\bottomrule
\end{tabular*}
\end{table}

\noindent\textbf{Matched increments and contributions.}
Let $T_{j,r}$ denote mean wall time over the 1,000 measured executions
of cumulative configuration $j$ in run $r$.
For $j=1,\ldots,4$, the paired increment is
\[
 \Delta T_{j,r}=T_{j,r}-T_{j-1,r},
\]
where $j=0$ is the baseline.
Means and sample standard deviations of these differences are computed
across the five matched runs.
Payload increments are computed analogously.
Table~\ref{tab:f-paired-increments} reports the resulting increments;
they characterize this fixed activation order, rather than isolated
component executions.
Stage shares divide mean increments by the final cumulative time or
payload, including the baseline.
The reported total SD is computed directly from the final configuration,
not by adding the stage SDs.

\begin{table}[!t]
\centering
\caption{Matched increments along the cumulative validation sequence.
Times are mean $\pm$ sample SD across five paired run differences;
Total is aggregated from the final cumulative configuration.}
\label{tab:f-paired-increments}
\footnotesize
\renewcommand{\arraystretch}{1.08}
\setlength{\tabcolsep}{3pt}
\begin{tabular*}{\columnwidth}{@{\hspace{3pt}}l@{\extracolsep{\fill}}cc@{\hspace{3pt}}}
\toprule
\textbf{Added checks}
& \shortstack{\textbf{Wall increment}\\\textbf{(ms)}}
& \shortstack{\textbf{Payload increment}\\\textbf{(KiB)}} \\
\midrule
Baseline & 0.00047 $\pm$ 0.00007 & 0.0215 \\
Lifecycle & 6.025 $\pm$ 0.029 & 15.348 \\
Handoff & 5.426 $\pm$ 0.057 & 7.099 \\
Provider & 2.909 $\pm$ 0.089 & 7.075 \\
Settlement & 2.343 $\pm$ 0.198 & 2.674 \\
\midrule
\textbf{Total} & \textbf{16.704 $\pm$ 0.109} & \textbf{32.217} \\
\bottomrule
\end{tabular*}
\end{table}

\noindent\textbf{Run-level variation.}
Table~\ref{tab:f-run-means} provides the run means underlying
Figure~\ref{fig:f-component-costs}(b).
For configuration $j$, define
\[
 \overline{T}_j=\frac{1}{5}\sum_{r=1}^{5}T_{j,r},
 \qquad
 s_j=\sqrt{\frac{1}{4}\sum_{r=1}^{5}
                 (T_{j,r}-\overline{T}_j)^2}.
\]
The plotted signed deviation and CV, expressed as percentages, are
\[
 d_{j,r}=100\left(\frac{T_{j,r}}{\overline{T}_j}-1\right),
 \qquad
 \mathrm{CV}_j=100\frac{s_j}{\overline{T}_j}.
\]
The zero line is each configuration's mean, and its shaded interval is
$[-\mathrm{CV}_j,+\mathrm{CV}_j]$ on the percentage axis.
Vertical offsets separate overlapping points and encode no additional
measurement.
These intervals describe one sample standard deviation of five run
means, rather than confidence intervals or per-execution variability.
Likewise, the time whiskers in Figure~\ref{fig:f-component-costs}(a)
describe cumulative run means; their SDs differ from those of the paired
increments in Table~\ref{tab:f-paired-increments}.
All statistics use unrounded observations.

\begin{table}[!t]
\centering
\caption{Run-mean cumulative validation time (ms).
Each column adds checks to the preceding configuration.}
\label{tab:f-run-means}
\footnotesize
\renewcommand{\arraystretch}{1.08}
\setlength{\tabcolsep}{3pt}
\begin{tabular*}{\columnwidth}{@{\hspace{3pt}}l@{\extracolsep{\fill}}cccc@{\hspace{3pt}}}
\toprule
\textbf{Run} & \textbf{Lifecycle} & \textbf{+ Handoff}
& \textbf{+ Provider} & \textbf{+ Settlement} \\
\midrule
1 & 6.027 & 11.480 & 14.318 & 16.721 \\
2 & 5.994 & 11.442 & 14.260 & 16.865 \\
3 & 6.042 & 11.472 & 14.379 & 16.618 \\
4 & 6.064 & 11.391 & 14.335 & 16.729 \\
5 & 6.002 & 11.473 & 14.512 & 16.589 \\
\midrule
CV (\%) & 0.478 & 0.321 & 0.661 & 0.652 \\
\bottomrule
\end{tabular*}
\end{table}
